\section{Positive multiplicities and complete canonical parent spaces}
\label{sec:peps_copy_parent_transport}

Fix a finite group $G$, a finite ordered simple graph, and a finite family
of matrix representations $D_i:G\to\operatorname{Mat}_{d_i}(\C)$.
Write $X=\coprod_i\{0,\ldots,d_i-1\}$ and
$X_m=\coprod_i(\{0,\ldots,d_i-1\}\times\{0,\ldots,m_i-1\})$.
An edge is written $f=(t,h)$; physical bond coordinates are ordered
head--tail. The construction extends the dimension multiplicities in
\cite[Section~7]{Schuch2010PEPS} to arbitrary positive integers $m_i$.
Irreducibility of $D_i$ is not required.

\subsection{Copy weights and arbitrary open boundaries}

\begin{definition}[Independent copy multiplicities]%
    \label{def:peps_copy_representation}
    \lean{TNLean.PEPS.blockMultiplicityRepresentation,
        TNLean.PEPS.blockMultiplicityRootWeight,
        TNLean.PEPS.blockMultiplicityRepresentation_self,
        TNLean.PEPS.blockMultiplicityRootWeight_self}
    \leanok
    Define $D(g)=\bigoplus_iD_i(g)$,
    $D^{[m]}(g)=\bigoplus_iD_i(g)\otimes I_{m_i}$, and
    $W_m=\bigoplus_i m_i^{1/4}I_{d_i}$.
    Taking $m_i=d_i$ gives dimension multiplicity restoration and the
    fourth-root dimension weight.
    Let $A_m$ be the actual averaging tensor for $D$ with virtual
    weight $W_m$, and let $B_m$ be the unweighted averaging tensor
    for $D^{[m]}$. Their site averages use the same factor $|G|^{-1}$.
\end{definition}

\begin{theorem}[One and two endpoint weights]%
    \label{thm:peps_copy_weight_products}
    \lean{TNLean.PEPS.blockMultiplicityRootWeight_mul,
        TNLean.PEPS.mul_blockMultiplicityRootWeight,
        TNLean.PEPS.blockMultiplicityRootWeight_commute,
        TNLean.PEPS.blockMultiplicityRootWeight_sq_mul}
    \leanok
    \uses{def:peps_copy_representation}
    For every $g\in G$,
    \begin{align}
        W_mD(g)=D(g)W_m & =\bigoplus_i m_i^{1/4}D_i(g),   \\
        W_m^2D(g)       & =\bigoplus_i\sqrt{m_i}\,D_i(g).
        \label{eq:peps_copy_weight_products}
    \end{align}
\end{theorem}
\begin{proof}\leanok
    Multiply within each diagonal block. Its scalar weight commutes
    with every matrix, and $(m_i^{1/4})^2=\sqrt{m_i}$.
\end{proof}

\begin{definition}[Regional copy transformations]%
    \label{def:peps_copy_regional_matrix}
    \lean{TNLean.PEPS.graphRegionCopyMatrix,
        TNLean.PEPS.graphRegionMultiplicityMatrix,
        TNLean.PEPS.numberedRegionCopyMatrix,
        TNLean.PEPS.numberedRegionMultiplicityMatrix}
    \leanok
    Let $F_m:\C^{X\times X}\to\C^{X_m\times X_m}$ be the supported
    normalized copy map. For $y=(i,a,c)$, write $\pi y=(i,a)$, and set
    $\alpha_m(y,z)=m_i^{-1/2}$ when $y,z$ have sector $i$ and the same
    copy label, and zero otherwise.
    Its one-ended map with fixed exterior output $z$ is
    $B_z(y,x)=\alpha_m(y,z)\delta_{\pi y,x}$.
    For a region $R$, let $E_R$ be its internal edges and
    $\partial E_R$ its crossing edges. Given
    $\gamma:\partial E_R\to X_m$, define
    \begin{align}
        M_{R,\gamma}
        =\bigotimes_{f\in\partial E_R}B_{\gamma_f}
        \otimes\bigotimes_{f\in E_R}F_m.
        \label{eq:peps_copy_regional_matrix}
    \end{align}
    The numbered regional matrix is obtained by regrouping the site
    labels into these crossing endpoints and internal head--tail pairs.
    The dimension-multiplicity versions are obtained by setting $m=d$.
\end{definition}

\begin{theorem}[Transport of internal and crossing factors]%
    \label{thm:peps_copy_open_factors}
    \lean{TNLean.PEPS.graphOpenInternalFactor_copy,
        TNLean.PEPS.graphOpenInternalFactor_copy_adjoint,
        TNLean.PEPS.graphOpenBoundaryFactor_copy,
        TNLean.PEPS.graphOpenBoundaryFactor_copy_adjoint,
        TNLean.PEPS.graphOpenInternalFactor_multiplicity,
        TNLean.PEPS.graphOpenInternalFactor_multiplicity_adjoint,
        TNLean.PEPS.graphOpenBoundaryFactor_multiplicity,
        TNLean.PEPS.graphOpenBoundaryFactor_multiplicity_adjoint}
    \leanok
    \uses{def:peps_copy_representation,def:peps_copy_regional_matrix}
    Assume $m_i>0$. Fix $q:R\to G$.
    Denote the internal bond factors of $A_m,B_m$ by $a_f(q),b_f(q)$,
    and the crossing factors with virtual label $\theta$ by
    $a_f(q;\theta),b_f(q;\theta)$.
    For an internal edge, $F_ma_f(q)=b_f(q)$ and
    $F_m^\dagger b_f(q)=a_f(q)$.
    For a crossing edge and exterior output $z\in X_m$,
    \begin{align}
        B_za_f(q;\theta)
         & =\sum_{y\in X_m}w_{y}\alpha_m(y,z)
        \delta_{\pi y,\theta}\,b_f(q;y),
        \label{eq:peps_copy_boundary_forward}  \\
        B_z^\dagger b_f(q;y)
         & =w_y^{-1}\alpha_m(y,z)a_f(q;\pi y),
        \label{eq:peps_copy_boundary_adjoint}
    \end{align}
    where $w_y=m_i^{1/4}$ for $y$ in sector $i$.
    These formulas hold for both crossing orientations; each crossing
    retains exactly one fourth-root weight.
\end{theorem}
\begin{proof}\leanok
    \uses{thm:peps_copy_weight_products,thm:peps_full_multiplicity_bond_map,
        thm:peps_parent_support_boundary_factors,
        thm:peps_parent_support_boundary_adjoints}
    The internal factor is $W_m^2D(q_hq_t^{-1})$.
    Equation~(\ref{eq:peps_copy_weight_products}) supplies
    $\sqrt{m_i}$ in each block, canceled by the normalized copy map.
    Its adjoint sums $m_i$ equal-copy terms with coefficient
    $m_i^{-1/2}$, recovering the same factor.
    At a crossing, use $W_mD(q_t^{-1})$ when the tail lies in $R$
    and $D(q_h)W_m$ when the head lies in $R$.
    The sparse coefficient of $B_z$ fixes the base coordinate and
    copy label; moving the remaining $w_y$ to the virtual boundary
    coefficient gives~(\ref{eq:peps_copy_boundary_forward}).
    Conjugating the same sparse coefficient and dividing by the
    nonzero $w_y$ gives~(\ref{eq:peps_copy_boundary_adjoint}).
\end{proof}

\begin{theorem}[Every virtual boundary coefficient is transported]%
    \label{thm:peps_copy_open_generators}
    \lean{TNLean.PEPS.graphRegionCopyMatrix_open,
        TNLean.PEPS.graphRegionCopyMatrix_adjoint_open,
        TNLean.PEPS.graphRegionMultiplicityMatrix_open,
        TNLean.PEPS.graphRegionMultiplicityMatrix_adjoint_open}
    \leanok
    \uses{def:peps_copy_regional_matrix}
    Assume $m_i>0$, and write $a_R(\theta),b_R(\eta)$ for the open
    contractions of $A_m,B_m$ in regional bond coordinates.
    For arbitrary virtual labels $\theta:\partial E_R\to X$ and
    $\eta:\partial E_R\to X_m$,
    \begin{align}
        M_{R,\gamma}a_R(\theta)
         & =\sum_\eta\left(\prod_{f\in\partial E_R}
        w_{\eta_f}\alpha_m(\eta_f,\gamma_f)
        \delta_{\pi\eta_f,\theta_f}\right)b_R(\eta),
        \label{eq:peps_copy_open_forward}           \\
        M_{R,\gamma}^\dagger b_R(\eta)
         & =\left(\prod_{f\in\partial E_R}
        w_{\eta_f}^{-1}\alpha_m(\eta_f,\gamma_f)\right)
        a_R(\pi\eta).
        \label{eq:peps_copy_open_adjoint}
    \end{align}
    Linear extension therefore transports arbitrary, possibly
    entangled, virtual boundary coefficients. These identities also
    give the dimension-multiplicity formulas when $m=d$.
\end{theorem}
\begin{proof}\leanok
    \uses{thm:peps_copy_open_factors,def:peps_parent_support_open_factors}
    Expand each open contraction as $|G|^{-|R|}$ times a sum over
    $q:R\to G$ of the product of its crossing and internal factors.
    Apply Theorem~\ref{thm:peps_copy_open_factors} to each factor.
    Distributing the crossing sums gives
    (\ref{eq:peps_copy_open_forward}); collecting their scalar
    coefficients gives~(\ref{eq:peps_copy_open_adjoint}).
    The transformation never restricts the original virtual boundary.
\end{proof}

\begin{theorem}[Both regional range inclusions]%
    \label{thm:peps_copy_local_ranges}
    \lean{TNLean.PEPS.graphRegionCopyMatrix_maps_openBondSpace,
        TNLean.PEPS.graphRegionCopyMatrix_adjoint_maps_openBondSpace,
        TNLean.PEPS.numberedRegionCopyMatrix_maps_regionGroundSpace,
        TNLean.PEPS.numberedRegionCopyMatrix_adjoint_maps_regionGroundSpace,
        TNLean.PEPS.graphRegionMultiplicityMatrix_maps_openBondSpace,
        TNLean.PEPS.graphRegionMultiplicityMatrix_adjoint_maps_openBondSpace,
        TNLean.PEPS.numberedRegionMultiplicityMatrix_maps_regionGroundSpace,
        TNLean.PEPS.numberedRegionMultiplicityMatrix_adjoint_maps_regionGroundSpace}
    \leanok
    If $m_i>0$, then for every region and every fixed exterior output
    $\gamma$, the regional matrix and its adjoint satisfy
    \begin{align}
        M_{R,\gamma}\mathcal G_R(A_m) & \subseteq\mathcal G_R(B_m),
                                      & M_{R,\gamma}^\dagger\mathcal G_R(B_m) & \subseteq\mathcal G_R(A_m).
        \label{eq:peps_copy_local_ranges}
    \end{align}
    The statements hold both in bond coordinates and after the
    supplied enumerations of the site physical spaces, including $m=d$.
\end{theorem}
\begin{proof}\leanok
    \uses{thm:peps_copy_open_generators,thm:peps_parent_support_open_coordinates}
    The vectors $a_R(\theta)$ and $b_R(\eta)$ span the respective
    full regional ranges. Equations~(\ref{eq:peps_copy_open_forward})
    and~(\ref{eq:peps_copy_open_adjoint}) send every generator into
    the required span. The coordinate permutations identify these
    spans with the original numbered contraction ranges.
\end{proof}

\subsection{Global supported inverses}

\begin{theorem}[All parent constraints are preserved]%
    \label{thm:peps_copy_parent_inclusions}
    \lean{TNLean.PEPS.graphNumberedCopyMatrix_maps_parent,
        TNLean.PEPS.graphNumberedCopyMatrix_adjoint_maps_parent,
        TNLean.PEPS.graphNumberedMultiplicityMatrix_maps_parent,
        TNLean.PEPS.graphNumberedMultiplicityMatrix_adjoint_maps_parent}
    \leanok
    Suppose $m_i>0$ and let $M_m$ be the actual product
    $\bigotimes_{f\in E}F_m$ in the supplied numbered site coordinates.
    For any family $\mathcal R=(R_i)_i$,
    \begin{align}
        M_m\mathcal P_{\mathcal R}(A_m) & \subseteq\mathcal P_{\mathcal R}(B_m),
                                        & M_m^\dagger\mathcal P_{\mathcal R}(B_m) & \subseteq\mathcal P_{\mathcal R}(A_m).
        \label{eq:peps_copy_parent_inclusions}
    \end{align}
    No coverage assumption on the regions is needed for these inclusions.
\end{theorem}
\begin{proof}\leanok
    \uses{thm:peps_copy_local_ranges,thm:peps_parent_support_copy_blocks,
        thm:peps_parent_support_block_criterion}
    Fix a region and the input and output exterior configurations.
    The corresponding block of $M_m$ is a scalar times
    $M_{R,\gamma}$: exterior edges give scalar factors, and crossing
    edges additionally enforce equality of the exterior base labels.
    Every such block carries the source regional range into the target
    range by~(\ref{eq:peps_copy_local_ranges}).
    Each target regional slice is a sum of these block images of
    source slices, and so satisfies the target condition.
    The adjoint block is the conjugate scalar times
    $M_{R,\gamma}^\dagger$, proving the reverse inclusion.
\end{proof}

\begin{theorem}[Exact inverses on the entire parent spaces]%
    \label{thm:peps_copy_parent_inverses}
    \lean{TNLean.PEPS.graphNumberedCopyMatrix_leftInverse_parent,
        TNLean.PEPS.graphNumberedCopyMatrix_rightInverse_parent,
        TNLean.PEPS.graphNumberedCopyMatrix_injOn_parent,
        TNLean.PEPS.graphNumberedMultiplicityMatrix_leftInverse_parent,
        TNLean.PEPS.graphNumberedMultiplicityMatrix_rightInverse_parent,
        TNLean.PEPS.graphNumberedMultiplicityMatrix_injOn_parent}
    \leanok
    Assume $m_i>0$ and every edge has both endpoints in some $R_i$.
    For $\psi\in\mathcal P_{\mathcal R}(A_m)$ and
    $\phi\in\mathcal P_{\mathcal R}(B_m)$,
    \begin{align}
        M_m^\dagger M_m\psi & =\psi,
                            & M_mM_m^\dagger\phi & =\phi.
        \label{eq:peps_copy_parent_inverses}
    \end{align}
    In particular, $M_m$ is injective on the source parent space.
    These identities concern every parent vector, including for $m=d$.
\end{theorem}
\begin{proof}\leanok
    \uses{thm:peps_copy_weight_products,thm:peps_parent_support_factor_ranges,
        thm:peps_parent_support_retractions}
    The source internal factors $\bigoplus_i\sqrt{m_i}D_i(g)$ lie
    in the matching-sector inclusion range. The target internal
    factors $\bigoplus_iD_i(g)\otimes I_{m_i}$ lie in $\operatorname{im}F_m$.
    Fixing crossing physical labels therefore puts every internal
    slice of a regional generator in the corresponding product range.
    Linear extension gives the same property on its regional range.
    For each edge, choose a parent region containing both endpoints;
    its regional constraint places every global one-bond slice in
    the required range. The product-range criterion then gives source
    matching-sector support and target product-image support.
    On these supports the product partial-isometry identities are
    exactly~(\ref{eq:peps_copy_parent_inverses}). Applying $M_m^\dagger$
    to $M_m\psi=M_m\psi'$ proves injectivity.
\end{proof}

\begin{theorem}[Canonical parent equivalence under copy restoration]%
    \label{thm:peps_copy_parent_equivalence}
    \lean{TNLean.PEPS.map_regionParentGroundSpace_copy_eq,
        TNLean.PEPS.canonicalCopyParentEquiv,
        TNLean.PEPS.finrank_regionParentGroundSpace_copy_eq,
        TNLean.PEPS.map_ker_regionParentHamiltonian_copy_eq,
        TNLean.PEPS.map_regionParentGroundSpace_multiplicity_eq,
        TNLean.PEPS.canonicalMultiplicityParentEquiv,
        TNLean.PEPS.finrank_regionParentGroundSpace_multiplicity_eq,
        TNLean.PEPS.map_ker_regionParentHamiltonian_multiplicity_eq}
    \leanok
    If $m_i>0$ and every edge is internal to some parent region,
    restriction of the actual product map gives
    \begin{align}
        M_m:\mathcal P_{\mathcal R}(A_m)
         & \xrightarrow{\ \cong\ }\mathcal P_{\mathcal R}(B_m),
         & \dim\mathcal P_{\mathcal R}(A_m)                     & =
        \dim\mathcal P_{\mathcal R}(B_m).
        \label{eq:peps_copy_parent_equivalence}
    \end{align}
    Its inverse is the restriction of $M_m^\dagger$.
    For any finite region family and positive interactions whose local
    kernels are these exact regional ranges, $M_m(\ker H)=\ker K$.
    Taking $m=d$ gives canonical dimension-multiplicity restoration.
\end{theorem}
\begin{proof}\leanok
    \uses{thm:peps_copy_parent_inclusions,thm:peps_copy_parent_inverses,
        thm:peps_region_parent_common_kernel}
    The inclusions in~(\ref{eq:peps_copy_parent_inclusions}) and both
    identities in~(\ref{eq:peps_copy_parent_inverses}) show that the
    restrictions are mutually inverse. The dimension assertion follows.
    Positivity identifies the Hamiltonian kernels with the common
    regional parent spaces, so the same map identifies those kernels.
\end{proof}

\subsection{Removing weights and changing multiplicities}

\begin{definition}[Physical copy filters]%
    \label{def:peps_copy_physical_filter}
    \lean{TNLean.PEPS.graphCopyWeight,
        TNLean.PEPS.graphCopyWeightPhysicalEquiv,
        TNLean.PEPS.toMatrix_graphCopyWeightPhysicalEquiv}
    \leanok
    For an incident physical configuration $\sigma$ at $v$, put
    $c_v(\sigma)=\prod_{f\in\operatorname{Inc}(v)}m_{i(\sigma_f)}^{1/4}$.
    When $m_i>0$, define the invertible diagonal physical filter
    $C_v=\operatorname{diag}(c_v)$ in the numbered site coordinates,
    and let $C=\bigotimes_vC_v$.
\end{definition}

\begin{theorem}[Copy weights are an invertible physical deformation]%
    \label{thm:peps_copy_physical_filter}
    \lean{TNLean.PEPS.graphDressedAveragingSite_copy_weight,
        TNLean.PEPS.graphCopyWeight_ne_zero,
        TNLean.PEPS.numberedGraphDressedAveragingTensor_copy_weight,
        TNLean.PEPS.map_regionParentGroundSpace_copy_weight,
        TNLean.PEPS.canonicalCopyWeightParentEquiv}
    \leanok
    \uses{def:peps_copy_physical_filter}
    Let $A_0$ be the unweighted averaging tensor for $D$.
    Every virtual column satisfies $(A_m)_v=C_v(A_0)_v$.
    If $m_i>0$, then every $c_v(\sigma)$ is nonzero and, for any region
    family without a coverage assumption,
    \begin{align}
        C\mathcal P_{\mathcal R}(A_0)=\mathcal P_{\mathcal R}(A_m),
        \qquad
        C:\mathcal P_{\mathcal R}(A_0)\cong\mathcal P_{\mathcal R}(A_m).
        \label{eq:peps_copy_filter_parent}
    \end{align}
\end{theorem}
\begin{proof}\leanok
    \uses{thm:peps_copy_weight_products,thm:peps_region_parent_ground_space_physical_transport}
    On each incident leg, the block-scalar weight can be read from the
    physical sector label. Taking the product of these weights gives
    $c_v(\sigma)$, independently of the virtual column.
    Positive $m_i$ make the diagonal matrices invertible.
    Physical deformation commutes with every open contraction.
    Applying $C$ and $C^{-1}$ to regional slices proves
    (\ref{eq:peps_copy_filter_parent}).
\end{proof}

\begin{definition}[The positive-multiplicity parent map]%
    \label{def:peps_copy_positive_parent_map}
    \lean{TNLean.PEPS.positiveMultiplicityParentMap,
        TNLean.PEPS.canonicalPositiveMultiplicityParentEquiv,
        TNLean.PEPS.canonicalMultiplicityChangeParentEquiv}
    \leanok
    \uses{def:peps_copy_physical_filter,thm:peps_copy_parent_equivalence}
    Assume all $m_i>0$, and that every edge is internal to some $R_i$.
    Define $T_m=M_mC$. Its restriction is the composite isomorphism
    from the unweighted single-copy parent space to that of $D^{[m]}$.
    For any two positive multiplicity families $m,n$, define their
    parent-space comparison by $T_nT_m^{-1}$, where the inverse is
    taken on the parent space.
\end{definition}

\begin{theorem}[Independence of positive copy multiplicities]%
    \label{thm:peps_copy_positive_parent}
    \lean{TNLean.PEPS.map_regionParentGroundSpace_positiveMultiplicity_eq,
        TNLean.PEPS.finrank_regionParentGroundSpace_positiveMultiplicity_eq,
        TNLean.PEPS.map_ker_regionParentHamiltonian_positiveMultiplicity_eq}
    \leanok
    \uses{def:peps_copy_positive_parent_map}
    For arbitrary representations $D_i$, positive $m_i$, supplied common
    incident enumerations, and internal coverage of every edge,
    \begin{align}
        T_m\mathcal P_{\mathcal R}(A_0) & =\mathcal P_{\mathcal R}(B_m),
                                        & \dim\mathcal P_{\mathcal R}(B_m) & =\dim\mathcal P_{\mathcal R}(A_0).
    \end{align}
    Consequently any two positive multiplicity families have isomorphic
    full canonical parent spaces. For a finite region family and
    arbitrary positive local interactions with the corresponding exact
    regional kernels, $T_m(\ker H)=\ker K$.
\end{theorem}
\begin{proof}\leanok
    \uses{thm:peps_copy_physical_filter,thm:peps_copy_parent_equivalence,
        thm:peps_region_parent_common_kernel}
    First apply $C$ using~(\ref{eq:peps_copy_filter_parent}), then apply
    $M_m$ using~(\ref{eq:peps_copy_parent_equivalence}).
    Both restrictions are isomorphisms, so their composition preserves
    dimensions. Reverse the comparison for $m$ and follow it by the
    comparison for $n$. The kernel statement again follows from
    positivity and the exact local-kernel hypotheses.
\end{proof}

% These results compare full kernels, not only chosen closed-state spans.
% Their hypotheses do not include a local-range equality, global-kernel
% equality, or a completeness claim for named torus closure vectors.
